\documentclass[11pt, a4paper]{article}
\usepackage[utf8]{inputenc}
\usepackage{amsmath, amssymb}
\usepackage{graphicx}
\usepackage{hyperref}
\usepackage{xcolor}
\usepackage{listings}
\lstset{
    language=Python,
    basicstyle=\ttfamily\small,
    keywordstyle=\color{blue}\bfseries,
    commentstyle=\color{gray}\itshape,
    stringstyle=\color{orange},
    numberstyle=\tiny\color{gray},
    numbers=left,
    stepnumber=1,
    numbersep=8pt,
    frame=single,
    breaklines=true,
    breakatwhitespace=false,
    showstringspaces=false,
    tabsize=4,
    captionpos=b
}
\usepackage{geometry}
\geometry{margin=1in}

\title{EE 325 Digital Signal Processing \\ Assignment 3: Discrete Time Convolution}
\author{PERERA J.D.T. \\ E/21/291}
\date{\today}

\begin{document}
\maketitle

\section{Objective}
Verify the calculation of linear convolution for discrete-time sequences using numerical tools.

\section{Q1: Linear Convolution}
\textbf{Sequences:}
\begin{itemize}
    \item $x[n] = [1, 2, 1]$
    \item $h[n] = [1, -1, 2]$
\end{itemize}

\textbf{Convolution $y[n] = x[n] * h[n]$:}
\begin{itemize}
    \item Length of $y$: $L_x + L_h - 1 = 3 + 3 - 1 = 5$.
    \item Result:
    \begin{itemize}
        \item $y[0] = 1$
        \item $y[1] = 1$
        \item $y[2] = 1$
        \item $y[3] = 3$
        \item $y[4] = 2$
    \end{itemize}
    \item \textbf{Verified Output:} $[1, 1, 1, 3, 2]$
\end{itemize}

\section{Q2: Cascade Connection of Systems}
\textbf{Systems:}
\begin{itemize}
    \item $h_1[n] = [1, 1]$
    \item $h_2[n] = [1, -1, 1]$
\end{itemize}

\textbf{Equivalent Impulse Response $h_{eq}[n]$:}
For cascade systems, $h_{eq}[n] = h_1[n] * h_2[n]$.
\begin{itemize}
    \item Length: $2 + 3 - 1 = 4$.
    \item Calculation:
    \begin{itemize}
        \item $h_{eq}[0] = 1$
        \item $h_{eq}[1] = 0$
        \item $h_{eq}[2] = 0$
        \item $h_{eq}[3] = 1$
    \end{itemize}
    \item \textbf{Verified Output:} $[1, 0, 0, 1]$
\end{itemize}

\section{Conclusion}
The computational results confirm that numerical convolution matches the analytical expectations for both single-stage and cascade LTI systems.

\appendix
\section{Python Source Code}
\textbf{File:} \texttt{convolution\_verification.py}
\begin{lstlisting}[caption={Discrete-Time Convolution Verification -- Assignment 3}]
import numpy as np
import matplotlib.pyplot as plt
import os

def plot_stem(n, y, title, filename):
    plt.figure(figsize=(8, 4))
    plt.stem(n, y, basefmt=" ")
    plt.title(title)
    plt.xlabel('n')
    plt.ylabel('Amplitude')
    plt.grid(True)
    plt.xticks(n)
    path = os.path.join('Convolution', 'plots', filename)
    os.makedirs(os.path.dirname(path), exist_ok=True)
    plt.savefig(path)
    print(f"Saved plot to {path}")
    plt.close()

def main():
    # Q1
    x = np.array([1, 2, 1])
    h = np.array([1, -1, 2])
    nx = np.arange(len(x)) # 0, 1, 2
    nh = np.arange(len(h)) # 0, 1, 2
    
    y1 = np.convolve(x, h)
    ny1 = np.arange(len(y1)) # 0, 1, 2, 3, 4
    
    print("Q1 Result y[n]:", y1)
    plot_stem(ny1, y1, 'Q1: Convolution y[n] = x[n] * h[n]', 'Q1_Output.png')

    # Q2
    h1 = np.array([1, 1])
    h2 = np.array([1, -1, 1])
    
    heq = np.convolve(h1, h2)
    nheq = np.arange(len(heq)) # 0 to 4
    
    print("Q2 Result heq[n]:", heq)
    plot_stem(nheq, heq, 'Q2: Cascade Response heq[n]', 'Q2_Output.png')

if __name__ == "__main__":
    main()
\end{lstlisting}

\end{document}
